\subsection{子群胚}

定义 5. —— 设 C 为范畴。C 的子范畴是 C 的满足下列条件的子箭图 D：

\begin{enumerate}
\def\labelenumi{(\roman{enumi})}
\item
  对 D 的每个顶点 a，\(e_a\) 是 D 的箭；
\item
  对 C 中属于 D 的每对可合成箭 \((f,g)\) ，积 fg 是 D 的箭；
\end{enumerate}

设 C 为范畴，D 为 C 的子范畴；赋有把 C 的合成法则限制到子集上而诱导的合成法则，D 是一个范畴。

C 的每个全子箭图都是 C 的子范畴。C 的一族子范畴的交是 C 的子范畴。

定义 6. —— 设 G 为群胚。G 的子群胚是 G 的这样一个子范畴 H：G 中属于 H 的每个箭的逆都是 H 的箭。

设 G 为群胚。G 的每个子群胚都是群胚。G 的每个全子箭图都是 G 的子群胚。G 的一族子群胚的交是 G 的子群胚。

设 H 为 G 的子箭图。以 H 为子箭图的 G 的诸子群胚之交称为由 H 生成的 G 的子群胚。其顶点集与 H 的顶点集相同，其轨道即 H 的连通分支。

例. —— 1) 设 X 为集合；以 X × X 表示 X 的配对群胚（II, 第 163 页, 例 2）。设 R 为 X 上的等价关系。以 X 为顶点集、以等价关系 R 的图为箭集的配对群胚 X × X 的子箭图，是 X × X 的一个子群胚。其轨道即关系 R 的等价类。

反之，X × X 中顶点集为 X 的每个子群胚都具有上述形式。

设 X 与 S 为集合，\(p: X \rightarrow S\) 为映射。映射 p 在 X 上定义一个等价关系（E, II, 第 41 页）。以 \(X \times_{S} X\) 表示由该等价关系定义的 \(X \times X\) 的子群胚，并称之为 \((X, p)\) 的配对群胚。

\begin{enumerate}
\def\labelenumi{\arabic{enumi})}
\setcounter{enumi}{1}
\item
  设 G 与 G' 为群胚，\(\varphi\) 与 \(\psi\) 为从 G 到 G' 的群胚态射。\(\varphi\) 与 \(\psi\) 的均衡子（参看 II, 第 153 页）是 G 的子群胚。
\item
  设 X 为集合，\(\Gamma\) 为群。以 \(X \times \Gamma \times X\) 表示这样的箭图：其顶点集为 X，其箭集为 \(X \times \Gamma \times X\) ，源映射与靶映射分别为 \(\mathrm{pr}_1\) 与 \(\mathrm{pr}_3\) 。赋有合成法则 \((x, \gamma, x') \cdot (x', \gamma', x'') = (x, \gamma \gamma', x'')\) 之后，它是一个群胚。它还是从群胚 \((\mathrm{X}\times\mathrm{X})\times\Gamma\) ——即配对群胚 \(X\times X\) 与群 \(\Gamma\) 相伴的群胚（II, 第 163 页, 例 1）之积——借助于双射 \((x,x',\gamma)\mapsto(x,\gamma,x')\) （它从 \(X\times X\times\Gamma\) 到 \(X\times\Gamma\times X\) ）经结构转移而得的群胚。
\end{enumerate}

设 \(m\colon X \times \Gamma \to X\) 为群 \(\Gamma\) 在 \(X\) 上的一个（右）作用。\(F\) （映射 \(m\) 的图像），也就是三元组 \((x, \gamma, x') \in X \times \Gamma \times X\) （满足 \(x' = x\gamma\) ）所成的集合，是唯一的子群胚 \(G\) （作为 \(X \times \Gamma \times X\) 的子群胚）的箭集。此群胚的轨道即 \(\Gamma\) 在 \(X\) 上的轨道；若 \(x \in X\) ，则 \(G\) 在 \(x\) 处的迷向群是形如 \((x, \gamma, x)\) 的元素所成的集合，其中 \(\gamma\) 取遍 \(x\) 在 \(\Gamma\) 中的稳定子。

反之，设 G 为群胚 \(X \times \Gamma \times X\) 的子群胚，且映射 \((\mathrm{pr}_{1}, \mathrm{pr}_{2})\) （从 \(\mathrm{Fl}(\mathrm{G})\) 到 \(X \times \Gamma\) ）是双射，则 G 具有上述形式。

\begin{enumerate}
\def\labelenumi{\arabic{enumi})}
\setcounter{enumi}{3}
\tightlist
\item
  设 G 为群胚，X 为集合，\(\varphi\colon\mathrm{X}\to\mathrm{Som}(\mathrm{G})\) 为映射。我们如下定义箭图 \(\mathrm{G}'\) ：集合 \(\mathrm{Som}(\mathrm{G}')\) 就是集合 X；对 X 中每一对元素 \((x,y)\) ，集合 \(\mathrm{Fl}_{x,y}(\mathrm{G}')\) 是三元组 \((x,f,y)\in\mathrm{X}\times\mathrm{Fl}(\mathrm{G})\times\mathrm{X}\) 的集合，其中 \(f\) 是 \(\mathrm{Fl}_{\varphi(x),\varphi(y)}(\mathrm{G})\) 的元素。设 \(x,y,z\) 为 X 的元素，\(f\in\mathrm{Fl}_{\varphi(x),\varphi(y)}(\mathrm{G})\) 与 \(g\in\mathrm{Fl}_{\varphi(y),\varphi(z)}(\mathrm{G})\) 为 G 的箭，并令 \((x,f,y)\cdot(y,g,z)=(x,fg,z)\) 。这就在箭图 \(\mathrm{G}'\) 中定义了一个合成法则，使之成为群胚。它称为群胚 G 经映射 \(\varphi\) 的逆像群胚，记作 \(\varphi^{*}(\mathrm{G})\) 。
\end{enumerate}

偶 \((\varphi,\psi)\) ——其中 \(\psi\colon\mathrm{Fl}(\varphi^{*}(\mathrm{G}))\to\mathrm{Fl}(\mathrm{G})\) 是由 \((x,f,y)\mapsto f\) 定义的映射——是从 \(\varphi^{*}(\mathrm{G})\) 到 G 的群胚态射，称为典范态射。

\begin{enumerate}
\def\labelenumi{\arabic{enumi})}
\setcounter{enumi}{4}
\tightlist
\item
  设 \(\varphi\colon\mathrm{G}\to\mathrm{G}^{\prime}\) 为群胚态射。设 H 为 G 的子箭图，其顶点集为 \(\mathrm{Som}(\mathrm{G})\) ，其箭集由 \(f\in\mathrm{Fl}(\mathrm{G})\) （满足 \(\varphi(f)\) 为 \(\mathrm{G}^{\prime}\) 的单位元）组成。对一切 \(a\in\mathrm{Som}(\mathrm{G})\) ，有 \(e_{a}\in\mathrm{H}_{a,a}\) 。对一切 \(f\in\mathrm{Fl}(\mathrm{H})\) ，有 \(\varphi(f^{-1})=\varphi(f)^{-1}\) ；因而 \(f^{-1}\in\mathrm{Fl}(\mathrm{H})\) 。此外，若 \(f\) 与 \(g\) 是 H 的可合成箭，则 \(\varphi(fg)=\varphi(f)\varphi(g)=e_{t(\varphi(f))}e_{o(\varphi(g))}=e_{t(\varphi(f))}\) ，因而 \(fg\) 是 H 的箭。这就证明了 H 是 G 的子群胚。它称为 \(\varphi\) 的核，记作 \(\mathrm{Ker}(\varphi)\) 。
\end{enumerate}

\pandocbounded{\includegraphics[keepaspectratio]{images/part002_cf856079.jpg}}

\begin{enumerate}
\def\labelenumi{\arabic{enumi})}
\setcounter{enumi}{5}
\tightlist
\item
  设 \(\varphi\colon\mathrm{G}\to\mathrm{G}^{\prime}\) 为群胚态射。箭图 \(\varphi(\mathrm{G})\) （其顶点集为 \(\varphi(\mathrm{Som}(\mathrm{G}))\) ，箭集为 \(\varphi(\mathrm{Fl}(\mathrm{G}))\) ）一般不是 \(\mathrm{G}^{\prime}\) 的子群胚。
\end{enumerate}

然而，若映射 Som(\(\varphi\)) 是单射，则它确为子群胚（II, 第 225 页, 习题 5）。

